% GENERATED FILE. Edit canonical lessons, then run npm run generate:books.
% canonical-source-hash: fnv1a64:b7f27222751fc771
% canonical-lessons: MW-C220-pakko, MW-C220-tyar, MW-C220-taalno, MW-C220-badalno, MW-C220-sodhno

\chapter{A Settled Plan}
\label{ch:mw-a-settled-plan}
\hlchaptermodality{220}

\begin{chapteropening}
\textbf{By the end of this chapter:} \emph{I can say my plan is settled, or that I will put it off or change it.}

Complete the last lesson of chapter 220: I can say my plan is settled, or that I will put it off or change it.
\end{chapteropening}

\section[\texorpdfstring{\mw{पक्को} (pakko)}{पक्को (pakko)}]{\mw{पक्को} (pakko) — settled; firm}

\label{lesson:MW-C220-pakko}



\begin{quote}
Before the new one: say the Marwadi for the Teej festival, then the Marwadi for in a hurry.
\end{quote}

\subsection*{You'll want to know: \mw{पक्को}}
\textbf{\mw{पक्को}} — \emph{pakko} — \textquotedblleft{}settled; firm\textquotedblright{}.

\begin{quote}
\mw{बात} \mw{पक्की} \mw{है}, \mw{हूं} \mw{काल} \mw{जासूं।}
\end{quote}

\begin{quote}
\emph{bāt pakkī hai, hū̃ kāl jāsū̃.}
\end{quote}

\begin{quote}
It's settled: I'm going tomorrow.
\end{quote}

\begin{cousinweb}
\textbf{\mw{पक्को}} (\emph{pakko}) goes back to Sanskrit \emph{pakva}, \textquotedblleft{}cooked, ripe\textquotedblright{}. Its doubled \emph{k} matches Hindi \emph{pakkā}, \textquotedblleft{}ripe, mature, real\textquotedblright{}, which CDIAL takes from Panjabi; the Old Marwari word CDIAL lists in the same entry is \emph{pāka}, \textquotedblleft{}ripe\textquotedblright{} (CDIAL 7621), the plain inherited shape. From \emph{ripe} it has come to mean \emph{firm}, \emph{settled}.
\end{cousinweb}

\begin{grammarlens}[title={What you've built}]
One more word for saying what comes next.
\end{grammarlens}

\subsection*{Your turn}
\begin{itemize}
\raggedright
  \item \emph{Say it:} \emph{pakko}
  \item \emph{Say it:} \emph{pakko}, once more
  \item \emph{Say it:} \emph{bāt pakkī hai, hū̃ kāl jāsū̃.} — It's settled: I'm going tomorrow.
  \item \emph{Recall:} say the Marwadi for the Teej festival, then the Marwadi for in a hurry, then say \emph{pakko} again
\end{itemize}

\begin{tcolorbox}[breakable,colback=teal!4,colframe=teal!35!black,title={Before you move on}]
What does \mw{पक्को} mean? (\textquotedblleft{}Settled; firm\textquotedblright{}.) Say it once more.
\end{tcolorbox}
\section[\texorpdfstring{\mw{त्यार} (tyār)}{त्यार (tyār)}]{\mw{त्यार} (tyār) — ready}

\label{lesson:MW-C220-tyar}



\begin{quote}
Before the new one: say the Marwadi for in a hurry, then the Marwadi for settled.
\end{quote}

\subsection*{You'll want to know: \mw{त्यार}}
\textbf{\mw{त्यार}} — \emph{tyār} — \textquotedblleft{}ready\textquotedblright{}.

\begin{quote}
\mw{हूं} \mw{अबार} \mw{त्यार} \mw{हूं।}
\end{quote}

\begin{quote}
\emph{hū̃ abār tyār hū̃.}
\end{quote}

\begin{quote}
I'm ready now.
\end{quote}

\begin{cousinweb}
\textbf{\mw{त्यार}} (\emph{tyār}) is usually traced to Persian \emph{taiyār}, \textquotedblleft{}ready\textquotedblright{}, which Persian took from Arabic; Marwadi has worn it down to one syllable.
\end{cousinweb}

\begin{grammarlens}[title={What you've built}]
One more word for saying what comes next.
\end{grammarlens}

\subsection*{Your turn}
\begin{itemize}
\raggedright
  \item \emph{Say it:} \emph{tyār}
  \item \emph{Say it:} \emph{tyār}, once more
  \item \emph{Say it:} \emph{hū̃ abār tyār hū̃.} — I'm ready now.
  \item \emph{Recall:} say the Marwadi for in a hurry, then the Marwadi for settled, then say \emph{tyār} again
\end{itemize}

\begin{tcolorbox}[breakable,colback=teal!4,colframe=teal!35!black,title={Before you move on}]
What does \mw{त्यार} mean? (\textquotedblleft{}Ready\textquotedblright{}.) Say it once more.
\end{tcolorbox}
\section[\texorpdfstring{\mw{टालणो} (ṭālṇo)}{टालणो (ṭālṇo)}]{\mw{टालणो} (ṭālṇo) — to put off; to avoid}

\label{lesson:MW-C220-taalno}



\begin{quote}
Before the new one: say the Marwadi for settled, then the Marwadi for ready.
\end{quote}

\subsection*{You'll want to know: \mw{टालणो}}
\textbf{\mw{टालणो}} — \emph{ṭālṇo} — \textquotedblleft{}to put off; to avoid\textquotedblright{}.

\begin{quote}
\mw{हूं} \mw{जातरा} \mw{कोनी} \mw{टालसूं।}
\end{quote}

\begin{quote}
\emph{hū̃ jātrā konī ṭālsū̃.}
\end{quote}

\begin{quote}
I won't put off the journey.
\end{quote}

Past (I …): \mw{म्हैं} \mw{जातरा} \mw{टाली}, \emph{mhaĩ jātrā ṭālī}.

\begin{cousinweb}
\textbf{\mw{टालणो}} (\emph{ṭālṇo}) is inherited from Sanskrit \emph{ṭalati}, \textquotedblleft{}moves aside\textquotedblright{}, through its causative \emph{ṭālayati}; CDIAL lists Marwari \emph{ṭālṇo}, \textquotedblleft{}to pass (time)\textquotedblright{}, and Hindi and Gujarati verbs for averting and avoiding under it (CDIAL 5450). To move a plan aside is to put it off.
\end{cousinweb}

\begin{grammarlens}[title={What you've built}]
One more word for saying what comes next.
\end{grammarlens}

\subsection*{Your turn}
\begin{itemize}
\raggedright
  \item \emph{Say it:} \emph{ṭālṇo}
  \item \emph{Say it:} \emph{ṭālṇo}, once more
  \item \emph{Say it:} \emph{hū̃ jātrā konī ṭālsū̃.} — I won't put off the journey.
  \item \emph{Recall:} say the Marwadi for settled, then the Marwadi for ready, then say \emph{ṭālṇo} again
\end{itemize}

\begin{tcolorbox}[breakable,colback=teal!4,colframe=teal!35!black,title={Before you move on}]
What does \mw{टालणो} mean? (\textquotedblleft{}To put off; to avoid\textquotedblright{}.) Say it once more.
\end{tcolorbox}
\section[\texorpdfstring{\mw{बदळणो} (badaḷṇo)}{बदळणो (badaḷṇo)}]{\mw{बदळणो} (badaḷṇo) — to change; to exchange}

\label{lesson:MW-C220-badalno}



\begin{quote}
Before the new one: say the Marwadi for ready, then the Marwadi for to put off.
\end{quote}

\subsection*{You'll want to know: \mw{बदळणो}}
\textbf{\mw{बदळणो}} — \emph{badaḷṇo} — \textquotedblleft{}to change; to exchange\textquotedblright{}.

\begin{quote}
\mw{हूं} \mw{काल} \mw{गेलो} \mw{बदळसूं।}
\end{quote}

\begin{quote}
\emph{hū̃ kāl gelo badaḷsū̃.}
\end{quote}

\begin{quote}
Tomorrow I'll change my route.
\end{quote}

Past (I …): \mw{म्हैं} \mw{गेलो} \mw{बदळ्यो}, \emph{mhaĩ gelo badaḷyo}.

\begin{cousinweb}
\textbf{\mw{बदळणो}} (\emph{badaḷṇo}) is usually traced to Persian \emph{badal}, \textquotedblleft{}an exchange, a substitute\textquotedblright{}, which Persian took from Arabic; Marwadi made a verb of it.
\end{cousinweb}

\begin{grammarlens}[title={What you've built}]
One more word for saying what comes next.
\end{grammarlens}

\subsection*{Your turn}
\begin{itemize}
\raggedright
  \item \emph{Say it:} \emph{badaḷṇo}
  \item \emph{Say it:} \emph{badaḷṇo}, once more
  \item \emph{Say it:} \emph{hū̃ kāl gelo badaḷsū̃.} — Tomorrow I'll change my route.
  \item \emph{Recall:} say the Marwadi for ready, then the Marwadi for to put off, then say \emph{badaḷṇo} again
\end{itemize}

\begin{tcolorbox}[breakable,colback=teal!4,colframe=teal!35!black,title={Before you move on}]
What does \mw{बदळणो} mean? (\textquotedblleft{}To change; to exchange\textquotedblright{}.) Say it once more.
\end{tcolorbox}
\section[\texorpdfstring{\mw{सोधणो} (sodhṇo)}{सोधणो (sodhṇo)}]{\mw{सोधणो} (sodhṇo) — to find out; to search out}

\label{lesson:MW-C220-sodhno}



\begin{quote}
Before the new one: say the Marwadi for to put off, then the Marwadi for to change.
\end{quote}

\subsection*{You'll want to know: \mw{सोधणो}}
\textbf{\mw{सोधणो}} — \emph{sodhṇo} — \textquotedblleft{}to find out; to search out\textquotedblright{}.

\begin{quote}
\mw{हूं} \mw{डूंगर} \mw{कनै} \mw{नुवो} \mw{गेलो} \mw{सोधसूं।}
\end{quote}

\begin{quote}
\emph{hū̃ ḍū̃gar kanai nuvo gelo sodhsū̃.}
\end{quote}

\begin{quote}
I'll search out a new path near the hill.
\end{quote}

Past (I …): \mw{म्हैं} \mw{गेलो} \mw{सोध्यो}, \emph{mhaĩ gelo sodhyo}.

\begin{cousinweb}
\textbf{\mw{सोधणो}} (\emph{sodhṇo}) is inherited from Sanskrit \emph{śoddhum}, \textquotedblleft{}to purify\textquotedblright{}; CDIAL lists the Old Marwari verb \emph{sodhaï}, \textquotedblleft{}finds out\textquotedblright{}, and a Gujarati verb \textquotedblleft{}to search\textquotedblright{} under it (CDIAL 12626).
\end{cousinweb}

\begin{grammarlens}[title={What you've built}]
One more word for saying what comes next.
\end{grammarlens}

\subsection*{Your turn}
\begin{itemize}
\raggedright
  \item \emph{Say it:} \emph{sodhṇo}
  \item \emph{Say it:} \emph{sodhṇo}, once more
  \item \emph{Say it:} \emph{hū̃ ḍū̃gar kanai nuvo gelo sodhsū̃.} — I'll search out a new path near the hill.
  \item \emph{Recall:} say the Marwadi for to put off, then the Marwadi for to change, then say \emph{sodhṇo} again
\end{itemize}

\begin{tcolorbox}[breakable,colback=teal!4,colframe=teal!35!black,title={Before you move on}]
What does \mw{सोधणो} mean? (\textquotedblleft{}To find out; to search out\textquotedblright{}.) Say it once more.
\end{tcolorbox}
